\documentclass[11pt,a4paper]{article}

\usepackage[margin=2.4cm]{geometry}
\usepackage[T1]{fontenc}
\usepackage[utf8]{inputenc}
\usepackage{lmodern}
\usepackage{setspace}
\usepackage{hyperref}
\usepackage{enumitem}

\title{LIF Networks in \textit{lif\_lova}:\\
\large A Musicological Description of Action, Form, and Harmonic Function}
\author{lif\_lova Documentation}
\date{\today}

\begin{document}
\maketitle
\onehalfspacing

\section*{Purpose and Scope}

This text is written for musicologists. Its aim is not to explain implementation details as
software engineering, but to describe how \textit{lif\_lova} transforms neural activity into
musical action, and how recognizable concepts from tonal and modal theory are applied in that
process.

In \textit{lif\_lova}, leaky integrate-and-fire (LIF) networks are not used as abstract visual
effects only. They function as an intermediary musical agent: they listen to audio-derived energy,
develop internal dynamic states, and externalize those states as tone and MIDI behavior.

\section*{General Musical Idea}

The system can be heard as a three-stage musical process:
\begin{enumerate}[label=\arabic*.]
    \item \textbf{Excitation}: incoming audio energy excites the network.
    \item \textbf{Organization}: network topology shapes how energy propagates, accumulates,
    and discharges.
    \item \textbf{Articulation}: the resulting activity is mapped to pitch, voicing,
    velocity, and temporal triggering.
\end{enumerate}

This resembles a historical distinction between \emph{material}, \emph{process}, and
\emph{surface}: signal material enters, form-generating process transforms it, and sounding surface
emerges as notes and timbre.

\section*{How LIF Action Becomes Musical Gesture}

The network continuously produces a distributed field of activation. A moving read position
(``scan phase'') samples a vertical slice of that field, yielding 16 ordered activity bins.
These bins serve as the immediate musical substrate.

From a musicological point of view, each bin acts like a voice region with two simultaneous roles:
\begin{itemize}
    \item \textbf{Spectral role}: in tone synthesis, bin intensity controls the contribution of a
    corresponding sine partial/oscillator.
    \item \textbf{Harmonic-role trigger}: in MIDI output, bin intensity and contour inform whether notes
    are activated, sustained, or released.
\end{itemize}

Thus, one and the same neural state is ``read'' in two idioms:
\begin{itemize}
    \item as \emph{continuous} amplitude in audio synthesis,
    \item and as \emph{discrete} event logic in MIDI harmony.
\end{itemize}

\section*{Application of Musical Theory Concepts}

\subsection*{1. Functional Harmony: Tonic, Subdominant, Dominant}

The system explicitly classifies activity into three functional categories:
\textbf{Tonic (T)}, \textbf{Subdominant (S)}, and \textbf{Dominant (D)}.

This is not a static chord lookup; it is a dynamic assignment influenced by:
\begin{itemize}
    \item bin position (register-like stratification),
    \item local energy prominence,
    \item contour (change relative to neighboring activity),
    \item and, in feedforward mode, an additional bias toward directed motion.
\end{itemize}

Musically, this means function is performed as behavior, not merely labeled as symbol.
The network can move from stability (T) toward preparation (S) and tension (D)
according to emergent activity patterns.

\subsection*{2. Modal Organization}

Pitch material is constrained by modal scales selected per scene:
Ionian, Dorian, Phrygian, Lydian, Mixolydian, Aeolian, and Locrian.

This creates a stable pitch grammar while preserving variability in voicing and rhythm.
In other words, the modal set defines permissible interval content, while the network defines
when and how that content is voiced.

\subsection*{3. Key and Register Framing}

The MIDI engine applies a transposition frame (key center by semitone) and a configurable
note-range window. Chord tones are octave-adjusted into this range.

This corresponds to two familiar compositional operations:
\begin{itemize}
    \item \textbf{Transpositional identity}: the same harmonic behavior in different keys.
    \item \textbf{Register orchestration}: keeping output in a chosen tessitura.
\end{itemize}

\subsection*{4. Style-Dependent Voicing Grammar}

Voicing strategy is style-selectable (Pop, Rock, Jazz, Blues, Percussion).
Each style modifies chord density and extension behavior while preserving
the underlying functional decision (T/S/D).

Seen musicologically, this is a layer of stylistic syntax over shared harmonic function:
\begin{itemize}
    \item \textbf{Pop}: compact consonant stacks with optional color extension.
    \item \textbf{Rock}: power-interval emphasis and stronger root-fifth profile.
    \item \textbf{Jazz}: denser tertian extension behavior.
    \item \textbf{Blues}: dominant color and blue-note inflection.
    \item \textbf{Percussion}: functional harmony suspended in favor of drum-note mapping.
\end{itemize}

\subsection*{5. Dynamics and Articulation}

Velocity is energy-shaped rather than fixed. Stronger and more prominent neural events yield
higher velocities, while weaker events remain softer. This creates phrasing that tracks
network intensity.

Two articulation regimes appear:
\begin{itemize}
    \item \textbf{Threshold sustain/release} (non-feedforward): activity crossing adaptive thresholds
    opens or closes note gates, producing held harmonic fields.
    \item \textbf{Pulse-gated triggering} (feedforward): short gate lengths plus cooldown intervals
    produce rhythmic, pointillistic behavior.
\end{itemize}

These regimes can be interpreted as contrasting temporal archetypes:
\emph{continuity} versus \emph{impulse sequence}.

\section*{Topology as Musical Form-Agent}

Network topology is musically consequential. It acts as a latent form model:
\begin{itemize}
    \item \textbf{Ring}: cyclical recurrence; periodic return and local continuity.
    \item \textbf{Fully connected}: high coupling; dense collective behavior and blended activity fields.
    \item \textbf{Feedforward}: directed flow; clearer attack chains and rhythmic propagation.
    \item \textbf{Sparse random}: punctuated, irregular link events; discontinuous accent structure.
    \item \textbf{Small-world}: local coherence with occasional long-range jumps;
    balance of stability and surprise.
\end{itemize}

For analysis, topology can be treated analogously to a formal constraint system:
it governs how quickly local events become global and how predictable recurrence is.

\section*{Dual Output: Sonification and Harmonic Agent}

The same LIF state produces two tightly related musical outputs:
\begin{enumerate}[label=\arabic*.]
    \item \textbf{Tone sonification}: 16-bin energy is rendered as an additive oscillator texture
    across a logarithmic frequency span.
    \item \textbf{MIDI harmony engine}: 16-bin energy drives functional classification,
    style-conditioned note selection, and event articulation.
\end{enumerate}

This duality is central: the listener may perceive the tone layer as a ``continuous shadow'' of
the same process that appears discretely in MIDI notes.

\section*{Interpretive Summary}

In \textit{lif\_lova}, LIF networks do not merely ``generate notes.''
They mediate between signal energy and music-theoretical organization.
Their action can be summarized as follows:
\begin{itemize}
    \item they convert audio excitation into structured dynamical behavior,
    \item they map that behavior into functional-harmonic categories,
    \item they realize those categories through modal, stylistic, and registral constraints,
    \item and they articulate the result through velocity and topology-dependent timing.
\end{itemize}

From a musicological perspective, this is a hybrid system where
computational dynamics and tonal-functional reasoning are co-compositional.

\section*{Keywords}

LIF network; computational musicology; functional harmony; modality; voicing;
algorithmic articulation; topology and form; sonification; MIDI harmony.

\end{document}